% GENERATED FILE. Edit canonical lessons, then run npm run generate:books.
% canonical-source-hash: fnv1a64:d3aa7aee6c2356b6
% canonical-lessons: TE-C320-uru, TE-S176-letter-uu, TE-C320-sutiga, TE-C320-ela-vellali, TE-R320-asking-the-way

\chapter{Asking the Way}
\label{ch:te-asking-the-way}
\hlchaptermodality{320}

\begin{chapteropening}
\textbf{By the end of this chapter:} \emph{I can ask where a place is and how to get there, follow directions with sūṭigā, eḍama vaipu and kuḍi vaipu, and talk about my ūru.}

Complete TE-R320-asking-the-way: ask the way, follow and give directions, and read the new letter \te{ఊ}.
\end{chapteropening}

\section[\texorpdfstring{\te{ఊరు} (ūru)}{ఊరు (ūru)}]{\te{ఊరు} (ūru) — a village or town; one's home place}

\label{lesson:TE-C320-uru}



\begin{quote}
Before the new one: say \textquotedblleft{}a village\textquotedblright{} (\emph{grāmaṁ}) and \textquotedblleft{}a town\textquotedblright{} (\emph{paṭṭaṇaṁ}). Then \textquotedblleft{}I must go\textquotedblright{}. (\emph{nēnu veḷḷāli}.)
\end{quote}

\subsection*{You'll want to know: \te{ఊరు}}
\textbf{\te{ఊరు}} — \emph{ūru} — \textquotedblleft{}a village or town; one's home place\textquotedblright{}.

\begin{quote}
\textbf{\te{మీది} \te{ఏ} \te{ఊరు}?}
\end{quote}

\begin{quote}
\emph{mīdi ē ūru?} — \textquotedblleft{}Which town are you from?\textquotedblright{}
\end{quote}

\begin{quote}
\textbf{\te{మా} \te{ఊరు} \te{చిన్నది}, \te{కానీ} \te{చాలా} \te{బాగుంటుంది}.}
\end{quote}

\begin{quote}
\emph{mā ūru cinnadi, kānī cālā bāguṇṭundi.} — \textquotedblleft{}My home town is small, but it is very nice.\textquotedblright{}
\end{quote}

\textbf{\te{మీది}} (\emph{mīdi}, \textquotedblleft{}yours\textquotedblright{}) and \textbf{\te{ఏ}} (\emph{ē}, \textquotedblleft{}which\textquotedblright{}): \textquotedblleft{}yours is which town?\textquotedblright{}, the usual way to ask where someone comes from.

\textbf{\te{ఊరు}} is wider than \textbf{\te{గ్రామం}} or \textbf{\te{పట్టణం}}: any place where people live, and most of all the one you come from. Before an ending its \textbf{-\te{ఉ}} turns to \textbf{-\te{ఇ}}: \textbf{\te{ఊరికి}} (\emph{ūriki}, \textquotedblleft{}to the town, home\textquotedblright{}), \textbf{\te{ఊరిలో}} (\emph{ūrilō}, \textquotedblleft{}in the town\textquotedblright{}). \textbf{\te{ఊరికి} \te{వెళ్తున్నాను}} on its own means \textquotedblleft{}I am going home, to my own place\textquotedblright{}. Many town names end in it, as \textbf{\te{నెల్లూరు}} does.

\begin{cousinweb}
\textbf{\te{ఊరు}} (\emph{ūru}) is native Dravidian, one word with Tamil and Malayalam \emph{ūr} and Kannada \emph{ūru}, all \textquotedblleft{}village, town\textquotedblright{} (DEDR 752). It is one of the words the whole family shares.
\end{cousinweb}

\subsection*{Your turn}
Say these aloud:

\begin{itemize}
\raggedright
  \item \emph{ūru}, with a long \emph{ū}
  \item \textquotedblleft{}Which town are you from?\textquotedblright{} — \emph{mīdi ē ūru?}
  \item \textquotedblleft{}My home town is small, but it is very nice\textquotedblright{} — \emph{mā ūru cinnadi, kānī cālā bāguṇṭundi}
  \item \textquotedblleft{}On holiday I go home to my town\textquotedblright{} — \emph{selavulō nēnu ūriki veḷtānu}
\end{itemize}

\begin{tcolorbox}[breakable,colback=teal!4,colframe=teal!35!black,title={Before you move on}]
What does \textbf{\te{ఊరు}} mean? (\textbf{A village or town; one's home place.}) \textquotedblleft{}To the town, home\textquotedblright{}? (\textbf{\te{ఊరికి}}.) How do you ask where someone comes from? (\textbf{\te{మీది} \te{ఏ} \te{ఊరు}?})
\end{tcolorbox}
\section[\texorpdfstring{\te{ఊ} (ū)}{ఊ (ū)}]{\te{ఊ} (ū) — the letter that opens \te{ఊరు}}

\label{lesson:TE-S176-letter-uu}



\begin{quote}
Before the new one: write \textbf{\te{ఉ}} and \textbf{\te{ూ}} once each, and name them. (Short \emph{u} at the start of a word; long \emph{ū} after a consonant, as in \textbf{\te{కూర}}.) Then write \textbf{\te{ఈ}} once.

Now say \textbf{\te{ఊరు}} (\emph{ūru}). This lesson takes its first letter out and puts it in your hand.
\end{quote}

\subsection*{Script you'll notice: \te{ఊ}}
\textbf{\te{ఊ}} — \emph{ū}, a long \emph{ū} at the start of a word.

After a consonant, long \emph{ū} is the sign \textbf{\te{ూ}}. At the start of a word it stands as a letter of its own, \textbf{\te{ఊ}}, and it pairs with the short \textbf{\te{ఉ}} you know.

In print \textbf{\te{ఊ}} is built on \textbf{\te{ఉ}}: the same body and bar, a top stroke that waves instead of a short tick, and a small curl standing on its right, like the curl beside \textbf{\te{ఈ}}.

\subsection*{Writing: \te{ఊ} — copy the printed shape}
Put your pen beside \textbf{\te{ఊ}} and copy it, slowly, and larger than it is printed. Write the body of \textbf{\te{ఉ}} the way you learned it, then copy the wave on top and the small curl on the right.

\begin{quote}
This book does not yet say where the pen starts on this shape, or which way it travels. That is taught with real variation from school to school, and it waits here until it can be written down from a source. Copying the printed shape needs no such source, and it puts the shape in your hand in the meantime.
\end{quote}

\subsection*{Your turn}
\begin{itemize}
\raggedright
  \item \emph{Find:} \te{ఊ} at the start of \textbf{\te{ఊరు}}
  \item \emph{Trace it:} \te{ఊ} three times, saying \emph{ū} as you finish each one
  \item \emph{Write it:} \te{ఉ}, then \te{ఊ}, from memory, naming each
  \item \emph{Write it:} \textbf{\te{ఊరు}}, then \textbf{\te{ఈ} \te{ఊరు}}, \textquotedblleft{}this town\textquotedblright{}
\end{itemize}

\begin{tcolorbox}[breakable,colback=teal!4,colframe=teal!35!black,title={Before you move on}]
Which letter is this — \te{ఊ}? (\textbf{ū}, long \emph{ū} at the start of a word.) Which word did it come from? (\textbf{\te{ఊరు}}, \emph{ūru}.) How is it different from \textbf{\te{ఉ}}? (\textbf{The wave on top and the curl on the right; the vowel is long.})
\end{tcolorbox}
\section[\texorpdfstring{\te{సూటిగా} (sūṭigā)}{సూటిగా (sūṭigā)}]{\te{సూటిగా} (sūṭigā) — straight on, directly}

\label{lesson:TE-C320-sutiga}



\begin{quote}
Before the new one: say \textquotedblleft{}slowly\textquotedblright{} (\emph{nemmadigā}) and \textquotedblleft{}quickly\textquotedblright{} (\emph{tvaragā}). Then say \textquotedblleft{}one's home place\textquotedblright{}. (\emph{ūru}.)
\end{quote}

\subsection*{You'll want to know: \te{సూటిగా}}
\textbf{\te{సూటిగా}} — \emph{sūṭigā} — \textquotedblleft{}straight on, directly\textquotedblright{}.

\begin{quote}
\textbf{\te{సూటిగా} \te{వెళ్ళండి}.}
\end{quote}

\begin{quote}
\emph{sūṭigā veḷḷaṇḍi.} — \textquotedblleft{}Go straight on.\textquotedblright{}
\end{quote}

\begin{quote}
\textbf{\te{అతను} \te{సూటిగా} \te{మాట్లాడతాడు}.}
\end{quote}

\begin{quote}
\emph{atanu sūṭigā māṭlāḍatāḍu.} — \textquotedblleft{}He speaks plainly, straight to the point.\textquotedblright{}
\end{quote}

Like \textbf{\te{నెమ్మదిగా}} and \textbf{\te{త్వరగా}}, it is a word with \textbf{-\te{గా}} on the end, and it answers \textquotedblleft{}how?\textquotedblright{}. On a road it is \textquotedblleft{}straight on\textquotedblright{}; about speech it is \textquotedblleft{}straight to the point\textquotedblright{}.

\begin{cousinweb}
\textbf{\te{సూటిగా}} (\emph{sūṭigā}) is native Dravidian: \textbf{\te{సూటి}}, \emph{sūṭi}, \textquotedblleft{}aim\textquotedblright{}, which DEDR 2658 lists with Tamil \emph{cuṭṭu}, \textquotedblleft{}to point out, to aim at\textquotedblright{}, and the adverb ending \textbf{-\te{గా}}. Going straight is going where you aim.
\end{cousinweb}

\subsection*{Your turn}
Say these aloud:

\begin{itemize}
\raggedright
  \item \emph{sūṭigā}
  \item \textquotedblleft{}Go straight on\textquotedblright{} — \emph{sūṭigā veḷḷaṇḍi}
  \item \textquotedblleft{}He speaks plainly, straight to the point\textquotedblright{} — \emph{atanu sūṭigā māṭlāḍatāḍu}
\end{itemize}

\begin{tcolorbox}[breakable,colback=teal!4,colframe=teal!35!black,title={Before you move on}]
What does \textbf{\te{సూటిగా}} mean? (\textbf{Straight on, directly.}) Which ending does it share with \textbf{\te{నెమ్మదిగా}}? (\textbf{-\te{గా}}.)
\end{tcolorbox}
\section[\texorpdfstring{\te{ఎలా} \te{వెళ్ళాలి}? (elā veḷḷāli?)}{ఎలా వెళ్ళాలి? (elā veḷḷāli?)}]{\te{ఎలా} \te{వెళ్ళాలి}? (elā veḷḷāli?) — asking the way}

\label{lesson:TE-C320-ela-vellali}



\begin{quote}
Say \textquotedblleft{}where?\textquotedblright{} (\emph{ekkaḍa}), \textquotedblleft{}how?\textquotedblright{} (\emph{elā}) and \textquotedblleft{}I must go\textquotedblright{} (\emph{nēnu veḷḷāli}). Then \textquotedblleft{}go straight on\textquotedblright{}. (\emph{sūṭigā veḷḷaṇḍi}.)
\end{quote}

\subsection*{You'll want to know: asking}
\begin{quote}
\textbf{\te{క్షమించండి}, \te{ఆసుపత్రి} \te{ఎక్కడ} \te{ఉంది}?}
\end{quote}

\begin{quote}
\emph{kṣamiñcaṇḍi, āsupatri ekkaḍa undi?} — \textquotedblleft{}Excuse me, where is the hospital?\textquotedblright{}
\end{quote}

\begin{quote}
\textbf{\te{ఆసుపత్రికి} \te{ఎలా} \te{వెళ్ళాలి}?}
\end{quote}

\begin{quote}
\emph{āsupatriki elā veḷḷāli?} — \textquotedblleft{}How do I get to the hospital?\textquotedblright{}
\end{quote}

The second is \textquotedblleft{}how must (I) go to …\textquotedblright{}, with the \textbf{-\te{ఆలి}} you learned for reasons.

\begin{quote}
\textbf{\te{ఇక్కడి} \te{నుంచి} \te{దూరమా}?}
\end{quote}

\begin{quote}
\emph{ikkaḍi nuñci dūramā?} — \textquotedblleft{}Is it far from here?\textquotedblright{}
\end{quote}

\textbf{-\te{ఆ}} at the end makes a yes-or-no question: \textbf{\te{దూరం}}, \textquotedblleft{}far\textquotedblright{} $\to$ \textbf{\te{దూరమా}?}, \textquotedblleft{}far?\textquotedblright{}. \textbf{\te{ఇక్కడి}} is the form \textbf{\te{ఇక్కడ}} takes before \textbf{\te{నుంచి}}.

\subsection*{You'll want to know: following the answer}
\begin{quote}
\textbf{\te{సూటిగా} \te{వెళ్ళండి}. \te{రెండవ} \te{మలుపులో} \te{ఎడమ} \te{వైపు} \te{తిరగండి}.}
\end{quote}

\begin{quote}
\emph{sūṭigā veḷḷaṇḍi. reṇḍava malupulō eḍama vaipu tiragaṇḍi.} — \textquotedblleft{}Go straight on. At the second turning, turn left.\textquotedblright{}
\end{quote}

\begin{quote}
\textbf{\te{ఆసుపత్రి} \te{కుడి} \te{వైపు} \te{ఉంది}. \te{దగ్గరే}.}
\end{quote}

\begin{quote}
\emph{āsupatri kuḍi vaipu undi. daggarē.} — \textquotedblleft{}The hospital is on the right. It is close.\textquotedblright{}
\end{quote}

\textbf{\te{వైపు}} (\emph{vaipu}, \textquotedblleft{}side\textquotedblright{}): \textbf{\te{ఎడమ} \te{వైపు}}, \textquotedblleft{}on the left, to the left\textquotedblright{}; \textbf{\te{కుడి} \te{వైపు}}, \textquotedblleft{}on the right\textquotedblright{}. \textbf{\te{తిరగండి}} (\emph{tiragaṇḍi}, \textquotedblleft{}turn\textquotedblright{}) is the polite command of a verb for turning, with the \textbf{-\te{అండి}} of \textbf{\te{క్షమించండి}}. \textbf{\te{దగ్గరే}} (\emph{daggarē}, \textquotedblleft{}close by\textquotedblright{}) is \textbf{\te{దగ్గర}} with an \textbf{-\te{ఏ}} that leans on it.

\subsection*{Your turn}
Say these aloud:

\begin{itemize}
\raggedright
  \item \textquotedblleft{}How do I get to the hospital?\textquotedblright{} — \emph{āsupatriki elā veḷḷāli?}
  \item \textquotedblleft{}Go straight on. At the second turning, turn left.\textquotedblright{} — \emph{sūṭigā veḷḷaṇḍi. reṇḍava malupulō eḍama vaipu tiragaṇḍi.}
  \item ask the way to the station, then give the way back: straight on, then right
\end{itemize}

\begin{tcolorbox}[breakable,colback=teal!4,colframe=teal!35!black,title={Before you move on}]
Ask \textquotedblleft{}how do I get there\textquotedblright{}, then \textquotedblleft{}is it far\textquotedblright{}. (\textbf{\te{ఎలా} \te{వెళ్ళాలి}}; \textbf{\te{దూరమా}}.) Then \textquotedblleft{}on the left\textquotedblright{} and \textquotedblleft{}turn\textquotedblright{}, politely. (\textbf{\te{ఎడమ} \te{వైపు}}; \textbf{\te{తిరగండి}}.)
\end{tcolorbox}
\section[(review)]{Review — asking the way}

\label{lesson:TE-R320-asking-the-way}



\begin{quote}
Say \textquotedblleft{}one's home place\textquotedblright{}, \textquotedblleft{}a vote\textquotedblright{} and \textquotedblleft{}in the meantime\textquotedblright{}. (\emph{ūru}, \emph{ōṭu}, \emph{īlōgā}.)

\begin{itemize}
\raggedright
  \item \emph{Write it:} the first letter of each, \textbf{\te{ఊ}}, \textbf{\te{ఓ}} and \textbf{\te{ఈ}}, once
\end{itemize}
\end{quote}

\subsection*{Your turn}
\begin{itemize}
\raggedright
  \item \emph{Say it:} Which town are you from? My home town is small.
  \item \emph{Say it:} Excuse me, how do I get to the station?
  \item \emph{Say it:} Go straight on. At the second turning, turn right. It is close.
  \item \emph{Say it:} It is far, so you had better take an auto.
  \item \emph{Write it:} \textbf{\te{ఊరు}}, \textbf{\te{ఓటు}}, \textbf{\te{ఈలోగా}}
\end{itemize}

\begin{tcolorbox}[breakable,colback=teal!4,colframe=teal!35!black,title={Before you move on}]
\textquotedblleft{}How do I get to the hospital?\textquotedblright{} (\textbf{\te{ఆసుపత్రికి} \te{ఎలా} \te{వెళ్ళాలి}?}) \textquotedblleft{}Straight on\textquotedblright{}? (\textbf{\te{సూటిగా}}.) \textquotedblleft{}On the right\textquotedblright{}? (\textbf{\te{కుడి} \te{వైపు}}.) \textquotedblleft{}One's home town\textquotedblright{}? (\textbf{\te{ఊరు}}, which opens with \textbf{\te{ఊ}}.)
\end{tcolorbox}
